\begin{homeworkProblem}[3]
  \begin{enumerate}[a)]
    \item Se $\lambda,\lambda_1,\lambda_2,\mu,\nu$ são medidas $\sigma$-finitas em $(X,\cali{A})$, mostre que
      \begin{enumerate}[1)]
        \item Se $\nu<<\lambda$ e $\lambda<<\mu$, então $\frac{d\nu}{d\mu}=\frac{d\nu}{d\lambda}\cdot\frac{d\lambda}{d\mu}$.
        \item Se $\lambda_1,\lambda_2<<\mu$, então $\frac{d(\lambda_1+\lambda_2)}{d\mu}=\frac{d\lambda_1}{d\mu}+\frac{d\lambda_2}{d\mu}$. 
      \end{enumerate}
    \item Definindo a variação total $|\lambda|$ de uma medida complexa $\lambda$ por $d|\lambda|=|f|d\mu$ (relação $d\lambda=fd\mu$ sempre que $\mu$ é positiva). Mostre que $|\lambda|$ está bem definida e valem as propriedadesdo item $a)$ também para medidas complexas absolutamente contínuas com respeito a uma medida positiva. 
  \end{enumerate}
  \begin{solution}
  \begin{enumerate}[a)]
    \item
    \begin{enumerate}[1)]
      \item Note que como $\nu<<\lambda$ e $\lambda<<\mu$, pelo teorema de Radon-Nikodym temos que existem $f,g\in M^{+}(X,\cali{A})$ tais que
        \begin{equation}\label{eq:1-3}
          \nu(A)=\int_{A}f\,d\lambda,\quad \lambda(A)=\int_{A}g\,d\mu,\quad\text{ para todo }A\in\cali{A}.
        \end{equation}
        Além disso, temos que o anterior implica que $\nu<<\mu$, pelo que pelo teorema de Radon-Nikodym temos que existe $h\in M^{+}(X,\cali{A})$ tal que
        \begin{align*}
          \nu(A)=\int_{A}h\,d\mu,\quad\text{ para todo }A\in\cali{A}. 
        \end{align*}
        Observe que, como $f\in M^{+}(X,\cali{A})$, temos que existe uma sequência monotona crescente $\{f_n\}_{n\in\mathbb{Z}^{+}}$ de funções simples tais que $f_n\to f$. Portanto, usando a equação \eqref{eq:1-3} e o Teorema da convergência monótona temos que
        \begin{align*}
          \nu(A)=\int_{A}f\,d\lambda=\lim_{n \to \infty}\int_{A}f_n\,d\lambda=\lim_{n \to \infty}\int_{A}\sum_{i=1}^{k}a_i\chi_{A_i}\,d\lambda=&\lim_{n \to \infty}\sum_{i=1}^{k}a_i \int_{A_i}\,d\lambda\\
          =&\lim_{n \to \infty}\sum_{i=1}^{k}a_i \lambda(A_i)\\
          =&\lim_{n \to \infty}\sum_{i=1}^{k}a_i \int_{A_i}g\,d\mu\\
          =&\lim_{n \to \infty}\int_{A}\sum_{i=1}^{k}a_i\chi_{A_i}g\,d\mu\\
          =&\lim_{n \to \infty}\int_{A}f_ng\,d\mu\\
          =&\int_{A}fg\,d\mu.
        \end{align*}
        Logo como a derivada de Radon-Nikodym é única q.t.p., temos que
        \begin{align*}
          \frac{d\nu}{d\mu}=h=fg=\frac{d\nu}{d\lambda}\cdot\frac{d\lambda}{d\mu}.
        \end{align*}
      \item Note que como $\lambda_1,\lambda_2<<\mu$ usando o Teorema de Radon-Nidokym temos que existem $f_1,f_2\in M^{+}(X,\cali{A})$ taes que
        \begin{align*}
          \lambda_1(A)=\int_{A}f_1\,d\mu,\quad \lambda_2(A)=\int_{A}f_2\,d\mu,\quad\text{ para todo }A\in\cali{A}.
        \end{align*}
        Além disso, temos que $\lambda_1+\lambda_2<<\mu$, pelo que pelo Teorema de Radon-Nikodym temos que existe $h\in M^{+}(X,\cali{A})$ tal que
        \begin{align*}
          (\lambda_1+\lambda_2)(A)=\int_{A}h\,d\mu,\quad\text{ para todo }A\in\cali{A}.
        \end{align*}
        Observe que como $f_1,f_2\in M^{+}(X,\cali{A})$ usando a linearidade da integral temos que 
        \begin{align*}
          (\lambda_1+\lambda_2)(A)=\lambda_1(A)+\lambda_2(A)=\int_{A}f_1\,d\mu+\int_{A}f_2\,d\mu=\int_{A}f_1+f_2\,d\mu.
        \end{align*}
        Logo, pela unicidade da derivada de Radon-Nikodym q.t.p., temos que
        \begin{align*}
          \frac{d(\lambda_1+\lambda_2)}{d\mu}=h=f_1+f_2=\frac{d\lambda_1}{d\mu}+\frac{d\lambda_2}{d\mu}.
        \end{align*}
    \end{enumerate}
    \item Sejam $\mu,\nu$ medidas positivas tais que $\lambda<<\mu$ e $\lambda<<\nu$. Logo, pelo Teorema de Radon-Nikodym (e como $\lambda$ é complexa) temos que existem $f\in L^{1}(\mu)$ e $g\in L^{1}(\nu)$ que satisfazen
      \begin{align*}
        \lambda(A)=\int_{A}f\,d\mu,\quad\lambda(A)=\int_{A}g\,d\nu,\quad\text{para todo }A\in\cali{A}.\\
      \end{align*}
      Logo definamos $|\lambda|$ e $\overline{|\lambda|}$ dadas por 
       \begin{align*}
        |\lambda|(A)=\int_{A}|f|\,d\mu,\quad\overline{|\lambda|}(A)=\int_{A}|g|\,d\nu,\quad\text{para todo }A\in\cali{A}.\\
      \end{align*}
      Vamos ver que $|\lambda|=\overline{|\lambda|}$.\\
      Definimos $\rho=\mu+\nu$. Note que $\mu<<\rho$ e $\nu<<\rho$, pelo que usando o Teorema de Radon-Nikodym, temos que existem $h_1,h_2\in M^{+}(X,\cali{A})$ tais que $d\mu=h_1\,d\rho$ e $d\nu=h_2\,d\rho$. Logo temos que usando a regra da cadeia mostrada no item $a)$ obtemos que 
      \begin{align*}
        \lambda(A)=\int_{A}f\,d\mu=\int_{A}fh_1\,d\rho.
      \end{align*}
      Similarmente
      \begin{align*}
        \lambda(A)=\int_{A}g\,d\nu=\int_{A}gh_2\,d\rho.
      \end{align*}
      Assim, pela unicidade da derivada de Radon-Nikodym, temos que $fh_1=gh_2$ $\rho$-q.t.p., o que implica que $|f|h_1=|g|h_2$ $\rho$-q.t.p.\\
      Segue-se que usando a regra da cadeia se cumpre que
      \begin{align*}
        |\lambda|(A)=\int_{A}|f|\,d\mu=\int_{A}|f|h_1\,d\rho=\int_{A}|g|h_2\,d\rho=\int_{A}|g|\,d\nu=\overline{|\lambda|}(A).        
      \end{align*}
      Pelo que podemos garantir que $|\lambda|=\overline{|\lambda|}$, i.e., $|\lambda|$ está bem definida.\\
      Agora, observe que se $\lambda_1$ e $\lambda_2$ são medidas complexas e $\mu$ é uma medida positiva tal que $\lambda_1<<\mu$ e $\lambda_2<<\mu$, temos que seguindo o mesmo razonamento do item $a)$ podemos verificar que
      \begin{align*}
        \frac{d(\lambda_1+\lambda_2)}{d\mu}=\frac{d\lambda_1}{d\mu}+\frac{d\lambda_2}{d\mu}.
      \end{align*}
\end{enumerate} 
  \end{solution}
\end{homeworkProblem}
